% Generated from extras.typ. Typst is authoritative; regeneration overwrites this file.
\documentclass{qlnotes-export}

\title{MATH 597: Measure Theory --- Supplementary Material}
\subtitle{Homework 0, problem-solving lectures, and practice problems}
\author{Qiulin Fan}
\date{Winter 2025}

\begin{document}
\mainmatter
\chapter{Use Egoroff and Hölder}\label{use-egoroff-and-huxf6lder}

Let \(\left\{ f_{n} \right\}\) be a sequence of functions in \(L^{p}\left( {\mathbb{R}}^{n} \right),1 < p < \infty\), which converge almost everywhere to a function \(f \in L^{p}\left( {\mathbb{R}}^{n} \right)\), and suppose that there is a constant \(M\) such that \(\parallel f_{n}\underset{p}{\parallel} \leq M\) for all \(n\). Show that for every \(g \in L^{q}\left( {\mathbb{R}}^{n} \right),q\) the conjugate of \(p\),

\[\int fg = \lim\limits_{n\rightarrow\infty}\int f_{n}g\]

Is the statement true for \(p = 1\) ? (Hint: you may want to use Egorov's Theorem.)
\end{document}
